\section{Limits and colimits of sheaves}
\label{section-limits-sheaves}

\noindent
Let $X$ be a topological space.
Let $\mathcal{I} \to \Sh(X)$, $i \mapsto \mathcal{F}_i$
be a diagram.
\begin{enumerate}
\item Both $\lim_i \mathcal{F}_i$ and $\colim_i \mathcal{F}_i$
exist.
\item The inclusion functor $i : \Sh(X) \to \textit{PSh}(X)$
commutes with limits. In other words, we may compute the limit
in the category of sheaves as the limit in the category of
presheaves. In particular, for any open $U \subset X$ we have
$$
(\lim_i \mathcal{F}_i)(U) = \lim_i \mathcal{F}_i(U).
$$
\item The inclusion functor $i : \Sh(X) \to \textit{PSh}(X)$
does not commute with colimits in general (not even
with finite colimits -- think surjections). The colimit is
computed as the sheafification of the colimit in the
category of presheaves:
$$
\colim_i \mathcal{F}_i =
\Big(U \mapsto \colim_i \mathcal{F}_i(U)\Big)^\#.
$$
\item Let $x \in X$ be a point. In general the stalk of
$\lim_i \mathcal{F}_i$ at $x$ is not equal to
the limit of the stalks. But if the index category is finite
then it is the case. In other words, the stalk functor is
left exact.
\item Let $x \in X$. We always have
$$
(\colim_i \mathcal{F}_i)_x = \colim_i \mathcal{F}_{i, x}.
$$
\item The sheafification functor
${}^\# : \textit{PSh}(X) \to \Sh(X)$ commutes with all
colimits, and with finite limits. But it does not commute
with all limits.
\end{enumerate}
The proofs are all easy. Here is an example of what is true for directed
colimits of sheaves.

\begin{lemma}
\label{lemma-directed-colimits-sections}
Let $X$ be a topological space. Let $I$ be a directed set.
Let $(\mathcal{F}_i, \varphi_{ii'})$ be a system of sheaves of sets
over $I$, see
Categories, Section \ref{categories-section-posets-limits}.
Let $U \subset X$ be an open subset.
Consider the canonical map
$$
\Psi :
\colim_i \mathcal{F}_i(U)
\longrightarrow
\left(\colim_i \mathcal{F}_i\right)(U)
$$
\begin{enumerate}
\item If all the transition maps are injective then
$\Psi$ is injective for any open $U$.
\item If $U$ is quasi-compact, then $\Psi$ is injective.
\item If $U$ is quasi-compact and all the transition maps are injective
then $\Psi$ is an isomorphism.
\item If $U$ has a cofinal system of open coverings
$\mathcal{U} : U = \bigcup_{j\in J} U_j$ with
$J$ finite and $U_j \cap U_{j'}$ quasi-compact
for all $j, j' \in J$, then $\Psi$ is bijective.
\end{enumerate}
\end{lemma}

\begin{proof}
Assume all the transition maps are injective. In this case the presheaf
$\mathcal{F}' : V \mapsto \colim_i \mathcal{F}_i(V)$ is
separated (see Definition \ref{definition-separated}).
By the discussion above we have
$(\mathcal{F}')^\# = \colim_i \mathcal{F}_i$.
By Lemma \ref{lemma-separated-presheaf-into-sheaf} we see that
$\mathcal{F}' \to (\mathcal{F}')^\#$ is injective. This proves (1).

\medskip\noindent
Assume $U$ is quasi-compact.
Suppose that $s \in \mathcal{F}_i(U)$ and
$s' \in \mathcal{F}_{i'}(U)$ give rise to elements on
the left hand side which have the same image under $\Psi$.
Since $U$ is quasi-compact this means there exists
a finite open covering $U = \bigcup_{j = 1, \ldots, m} U_j$
and for each $j$ an index $i_j \in I$, $i_j \geq i$, $i_j \geq i'$
such that $\varphi_{ii_j}(s)$ and $\varphi_{i'i_j}(s')$ have the
same restriction to $U_j$.
Let $i''\in I$ be $\geq$ than all of the $i_j$.
We conclude that $\varphi_{ii''}(s)$ and $\varphi_{i'i''}(s')$
agree on the opens $U_j$ for all $j$ and hence that
$\varphi_{ii''}(s) = \varphi_{i'i''}(s')$. This proves (2).

\medskip\noindent
Assume $U$ is quasi-compact and all transition maps injective.
Let $s$ be an element of the target of $\Psi$.
Since $U$ is quasi-compact
there exists a finite open covering $U = \bigcup_{j = 1, \ldots, m} U_j$,
for each $j$ an index $i_j \in I$ and $s_j \in \mathcal{F}_{i_j}(U_j)$
such that $s|_{U_j}$ comes from $s_j$ for all $j$.
Pick $i \in I$ which is $\geq$ than all of the $i_j$.
By (1) the sections $\varphi_{i_ji}(s_j)$ agree over the
overlaps $U_j \cap U_{j'}$. Hence they glue to a section
$s' \in \mathcal{F}_i(U)$ which maps to $s$ under $\Psi$.
This proves (3).

\medskip\noindent
Assume the hypothesis of (4). In particular we see that
$U$ is quasi-compact and hence by (2) we have injectivity of $\Psi$.
Let $s$ be an element of the target of $\Psi$.
By assumption there  exists a finite open covering
$U = \bigcup_{j = 1, \ldots, m} U_j$, with $U_j \cap U_{j'}$
quasi-compact for all $j, j' \in J$ and
for each $j$ an index $i_j \in I$ and $s_j \in \mathcal{F}_{i_j}(U_j)$
such that $s|_{U_j}$ is the image of $s_j$ for all $j$.
Since $U_j \cap U_{j'}$ is quasi-compact we can apply (2)
and we see that there exists an $i_{jj'} \in I$,
$i_{jj'} \geq i_j$, $i_{jj'} \geq i_{j'}$ such that
$\varphi_{i_ji_{jj'}}(s_j)$ and $\varphi_{i_{j'}i_{jj'}}(s_{j'})$
agree over $U_j \cap U_{j'}$. Choose an index $i \in I$
which is bigger or equal than all the $i_{jj'}$. Then we see that
the sections $\varphi_{i_ji}(s_j)$ of $\mathcal{F}_i$ glue
to a section of $\mathcal{F}_i$ over $U$. This section is mapped
to the element $s$ as desired.
\end{proof}

\begin{example}
\label{example-conditions-needed-colimit}
Let $X = \{s_1, s_2, \xi_1, \xi_2, \xi_3, \ldots\}$ as a set.
Declare a subset $U \subset X$ to be open
if $s_1 \in U$ or $s_2 \in U$ implies $U$ contains all of the
$\xi_i$. Let $U_n = \{\xi_n, \xi_{n + 1}, \ldots\}$, and
let $j_n : U_n \to X$ be the inclusion map.
Set $\mathcal{F}_n = j_{n, *}\underline{\mathbf{Z}}$.
There are transition maps $\mathcal{F}_n \to \mathcal{F}_{n + 1}$.
Let $\mathcal{F} = \colim \mathcal{F}_n$.
Note that $\mathcal{F}_{n, \xi_m} = 0$ if $m < n$
because $\{\xi_m\}$ is an open subset of $X$ which misses $U_n$.
Hence we see that $\mathcal{F}_{\xi_n} = 0$ for all $n$.
On the other hand the stalk $\mathcal{F}_{s_i}$, $i = 1, 2$
is the colimit
$$
M = \colim_n \prod\nolimits_{m \geq n} \mathbf{Z}
$$
which is not zero. We conclude that the sheaf
$\mathcal{F}$ is the direct sum of the
skyscraper sheaves with value $M$ at the closed points
$s_1$ and $s_2$. Hence $\Gamma(X, \mathcal{F}) = M \oplus M$.
On the other hand, the reader can verify that
$\colim_n \Gamma(X, \mathcal{F}_n) = M$.
Hence some condition is necessary in part (4)
of Lemma \ref{lemma-directed-colimits-sections} above.
\end{example}

\noindent
There is a version of the previous lemma dealing with sheaves on a diagram
of spectral spaces. To state it we introduce some notation.
Let $\mathcal{I}$ be a cofiltered index category.
Let $i \mapsto X_i$ be a diagram of spectral spaces over $\mathcal{I}$
such that for $a : j \to i$ in $\mathcal{I}$ the corresponding map
$f_a : X_j \to X_i$ is spectral. Set $X = \lim X_i$ and denote
$p_i : X \to X_i$ the projection.

\begin{lemma}
\label{lemma-compute-pullback-to-limit}
In the situation described above, let $i \in \Ob(\mathcal{I})$ and let
$\mathcal{G}$ be a sheaf on $X_i$. For $U_i \subset X_i$
quasi-compact open we have
$$
p_i^{-1}\mathcal{G}(p_i^{-1}(U_i)) =
\colim_{a : j \to i} f_a^{-1}\mathcal{G}(f_a^{-1}(U_i))
$$
\end{lemma}

\begin{proof}
Let us prove the canonical map
$\colim_{a : j \to i} f_a^{-1}\mathcal{G}(f_a^{-1}(U_i)) \to
p_i^{-1}\mathcal{G}(p_i^{-1}(U_i))$ is injective.
Let $s, s'$ be sections of $f_a^{-1}\mathcal{G}$
over $f_a^{-1}(U_i)$ for some $a : j \to i$. For
$b : k \to j$ let $Z_k \subset f_{a \circ b}^{-1}(U_i)$
be the closed subset of points $x$ such that the image of
$s$ and $s'$ in the stalk $(f_{a \circ b}^{-1}\mathcal{G})_x$
are different. If $Z_k$ is nonempty for all $b : k \to j$, then by
Topology, Lemma \ref{topology-lemma-inverse-limit-spectral-spaces-nonempty}
we see that $\lim_{b : k \to j} Z_k$ is nonempty too.
Then for $x \in \lim_{b : k \to j} Z_k \subset X$
(observe that $\mathcal{I}/j \to \mathcal{I}$ is initial)
we see that the image of $s$ and $s'$ in the stalk of
$p_i^{-1}\mathcal{G}$ at $x$ are different too since
$(p_i^{-1}\mathcal{G})_x = (f_{b \circ a}^{-1}\mathcal{G})_{p_k(x)}$
for all $b : k \to j$ as above. Thus if the images of $s$ and $s'$
in $p_i^{-1}\mathcal{G}(p_i^{-1}(U_i))$ are the same, then $Z_k$
is empty for some $b : k \to j$. This proves injectivity.

\medskip\noindent
Surjectivity. Let $s$ be a section of $p_i^{-1}\mathcal{G}$ over
$p_i^{-1}(U_i)$. By
Topology, Lemma \ref{topology-lemma-directed-inverse-limit-spectral-spaces}
the set $p_i^{-1}(U_i)$ is a quasi-compact open of the spectral space $X$.
By construction of the pullback sheaf, we can find an open covering
$p_i^{-1}(U_i) = \bigcup_{l \in L} W_l$,
opens $V_{l, i} \subset X_i$,
sections $s_{l, i} \in \mathcal{G}(V_{l, i})$
such that $p_i(W_l) \subset V_{l, i}$ and $p_i^{-1}s_{l, i}|_{W_l} = s|_{W_l}$.
Because $X$ and $X_i$ are spectral and $p_i^{-1}(U_i)$ is quasi-compact open,
we may assume $L$ is finite and $W_l$ and $V_{l, i}$ quasi-compact open
for all $l$. Then we can apply
Topology, Lemma \ref{topology-lemma-descend-opens}
to find $a : j \to i$ and open covering
$f_a^{-1}(U_i) = \bigcup_{l \in L} W_{l, j}$ by quasi-compact
opens whose pullback to $X$ is the covering
$p_i^{-1}(U_i) = \bigcup_{l \in L} W_l$
and such that moreover $W_{l, j} \subset f_a^{-1}(V_{l, i})$.
Write $s_{l, j}$ the restriction of the pullback of $s_{l, i}$
by $f_a$ to $W_{l, j}$. Then we see that $s_{l, j}$ and $s_{l', j}$
restrict to elements of $(f_a^{-1}\mathcal{G})(W_{l, j} \cap W_{l', j})$
which pullback to the same element $(p_i^{-1}\mathcal{G})(W_l \cap W_{l'})$,
namely, the restriction of $s$. Hence by injectivity, we can find
$b : k \to j$ such that the sections $f_b^{-1}s_{l, j}$
glue to a section over $f_{a \circ b}^{-1}(U_i)$ as desired.
\end{proof}

\noindent
Next, in addition to the cofiltered system $X_i$ of spectral spaces,
assume given
\begin{enumerate}
\item a sheaf $\mathcal{F}_i$ on $X_i$ for all $i \in \Ob(\mathcal{I})$,
\item for $a : j \to i$ an $f_a$-map
$\varphi_a : \mathcal{F}_i \to \mathcal{F}_j$
\end{enumerate}
such that $\varphi_c = \varphi_b \circ \varphi_a$
whenever $c = a \circ b$. Set $\mathcal{F} = \colim p_i^{-1}\mathcal{F}_i$
on $X$.

\begin{lemma}
\label{lemma-descend-opens}
In the situation described above, let $i \in \Ob(\mathcal{I})$ and let
$U_i \subset X_i$ be a quasi-compact open. Then
$$
\colim_{a : j \to i} \mathcal{F}_j(f_a^{-1}(U_i)) = \mathcal{F}(p_i^{-1}(U_i))
$$
\end{lemma}

\begin{proof}
Recall that $p_i^{-1}(U_i)$ is a quasi-compact open of the spectral space
$X$, see
Topology, Lemma \ref{topology-lemma-directed-inverse-limit-spectral-spaces}.
Hence Lemma \ref{lemma-directed-colimits-sections} applies and we have
$$
\mathcal{F}(p_i^{-1}(U_i)) =
\colim_{a : j \to i} p_j^{-1}\mathcal{F}_j(p_i^{-1}(U_i)).
$$
A formal argument shows that
$$
\colim_{a : j \to i} \mathcal{F}_j(f_a^{-1}(U_i)) =
\colim_{a : j \to i} \colim_{b : k \to j}
f_b^{-1}\mathcal{F}_j(f_{a \circ b}^{-1}(U_i))
$$
Thus it suffices to show that
$$
p_j^{-1}\mathcal{F}_j(p_i^{-1}(U_i)) =
\colim_{b : k \to j} f_b^{-1}\mathcal{F}_j(f_{a \circ b}^{-1}(U_i))
$$
This is Lemma \ref{lemma-compute-pullback-to-limit}
applied to $\mathcal{F}_j$ and the quasi-compact open $f_a^{-1}(U_i)$.
\end{proof}

















